Small models are becoming dense with useful intelligence. A few years ago, an edge gateway could realistically host lightweight classifiers, rules, or distilled neural networks; today, a 1--3B parameter small language model can support classification, structured interaction, local explanation, and task-specific reasoning. This changes the shape of edge intelligence. The question is no longer only how to compress a cloud model for inference, but how to let a compact local model become a durable operational assistant that can improve over time without exporting sensitive data.

Industrial Internet of Things (IIoT) security is a natural test case for this shift. Edge gateways collect packet statistics, protocol traces, sensor activity, logs, and system events that reveal sensitive operational behavior. Sending these records, or even rich prompts derived from them, to an external API creates privacy, compliance, latency, cost, and availability concerns. At the same time, a single gateway only sees a partial slice of normal behavior and attack patterns. A useful edge security assistant should therefore evolve collaboratively across sites while keeping raw telemetry local.

Federated learning provides the collaborative learning path, but it is not a complete privacy solution: individual updates can leak information, and an honest-but-curious server can inspect uploads before aggregation. Client-side release DP gives a stronger and auditable guarantee because each client clips and noises its update before transmission. The price is utility loss. For language-model adapters, this price is amplified by communication limits, non-IID gateway data, and the large number of trainable coordinates that may receive DP noise.

We instantiate this setting with Qwen3.5-2B~\citep{qwen35modelcard}, a recent open-weight model in the 0.5--3B edge-relevant range. This keeps the paper grounded in a concrete edge-small-model regime rather than a generic LLM setting. Although Qwen3.5-2B supports multimodal usage, our benchmark uses only the text/language pathway with telemetry rendered as compact textual prompts. The goal is not to show that a small language model dominates all tabular IDS classifiers; classical models remain important baselines. The goal is to study whether a compact edge assistant can be privately federated under realistic privacy and communication budgets.

Parameter-efficient fine-tuning methods such as LoRA~\citep{hu2021lora} reduce trainable parameters, but federated averaging of adapter factors can be suboptimal because the product parameterization is not uniquely identified. Recent fixed-basis and square-core approaches, including Fed-SB~\citep{singhal2025fedsb}, address this issue by training a small matrix between fixed adapter factors. \method does not claim that square-core training or exact core averaging is new. Instead, it asks a sharper question for edge intelligence: can public spectral structure from the released backbone be used to define, budget, privatize, and aggregate the minimal coordinates needed for a small edge model to keep learning?

\method answers this question with a public-spectral, budget-aware federated core adaptation design. For each selected linear layer of a frozen small language model, \method derives orthonormal bases from the public backbone weights, trains only a small layer-wise core matrix, clips and noises client uploads in core coordinates, and aggregates the privatized updates on the server. The method is designed around four principles: private-data-free basis construction, public spectrum-driven rank and noise allocation, client-level DP accounting, and local residual personalization for non-IID gateways.

The paper makes the following contributions:
\begin{itemize}[leftmargin=*]
  \item We formulate private federated evolution of edge small language models as a concrete IIoT security problem, emphasizing local data sovereignty, communication cost, client-level DP, rare attack recall, and reproducible evaluation.
  \item We introduce a public-spectral core adaptation mechanism that constructs all adapter bases from released backbone weights and trains only small core matrices, giving unique coordinates and exact aggregation inside a public subspace.
  \item We propose a DP utility-recovery stack for edge SLM adaptation: public spectrum-driven rank allocation, layer-adaptive clipping/noise allocation, local residual cores, and server-side post-processing.
  \item We define a Qwen3.5-2B Edge-IIoT instruction IDS protocol with fixed splits, simulated IID and Dirichlet non-IID clients, leakage diagnostics, raw-file robustness checks, and fair comparisons under matched privacy, communication, and parameter budgets.
\end{itemize}

This framing deliberately avoids claiming that a square core or exact core averaging is unique. The intended novelty is the combination of public-backbone spectral construction, DP utility recovery, edge-budget allocation, and an IIoT-specific protocol for continuously improving a compact private edge assistant.
